\documentclass[../DoAn.tex]{subfiles}
\begin{document}

\begin{center}
    \Large{\textbf{TÓM TẮT NỘI DUNG ĐỒ ÁN}}\\
\end{center}
\vspace{1cm}

Nhu cầu quan sát và phân tích mô phỏng giao thông ngày càng tăng, tuy nhiên
giao diện 2D truyền thống của \gls{SUMO} còn hạn chế về khả năng trực quan hóa
không gian, theo dõi đồng thời nhiều đối tượng và trình bày kịch bản cho người
không chuyên. Một số giải pháp hiển thị 3D đã tồn tại nhưng thường phụ thuộc
vào bản đồ tĩnh dựng sẵn, khó tái sử dụng cho nhiều kịch bản và ít hỗ trợ tương
tác trực tiếp với phương tiện. Trong đồ án này, em lựa chọn hướng tiếp cận kết
nối chuỗi công cụ \gls{SUMO} --- \gls{TraCI}/Python --- Unity: \gls{SUMO} mô
phỏng giao thông vi mô và sinh dữ liệu, Python dùng \gls{TraCI} thu thập trạng
thái theo từng bước thời gian và truyền sang Unity qua TCP/IP, còn Unity đảm
nhận dựng cảnh và hiển thị 3D theo thời gian thực. Hướng tiếp cận này được chọn
vì tận dụng được năng lực mô phỏng mạnh của \gls{SUMO} đồng thời khai thác
pipeline dựng hình, vật lý và tương tác của Unity. Trên cơ sở đó, hệ thống tự
động dựng mạng đường từ dữ liệu mô phỏng, hiển thị phương tiện, đèn tín hiệu và
người đi bộ, hỗ trợ hai chế độ vận hành là thời gian thực và phát lại, cùng các
thao tác điều khiển camera, tạm dừng/tiếp tục và điều chỉnh tốc độ mô phỏng.
Đóng góp nổi bật của đồ án là cơ chế cho phép người dùng chiếm quyền điều khiển
một phương tiện và lái thủ công bằng vật lý, xử lý va chạm để sinh ``xác xe'' gây
ùn ứ, và trả quyền điều khiển để xóa xe khỏi mô phỏng ngay lập tức; bên cạnh đó là cơ chế lưu
trữ phiên mô phỏng hợp nhất phục vụ truy vết và phát lại. Hệ thống còn được tối
ưu luồng truyền dữ liệu (gom truy vấn \gls{TraCI}, nén dữ liệu trạng thái, giảm
độ trễ gói tin) kèm cơ chế đo độ trễ đầu--cuối, và dựng thêm khối công trình từ
bản đồ OpenStreetMap để tăng tính trực quan. Kết quả thực nghiệm
trên các kịch bản sinh từ mê cung và từ bản đồ OpenStreetMap cho thấy hệ thống
dựng cảnh đúng dữ liệu nguồn, tái hiện chuyển động ổn định và đáp ứng nhu cầu
quan sát, phân tích trực quan.

\begin{flushright}
Sinh viên thực hiện\\
\begin{tabular}{@{}c@{}}
\textit{(Ký và ghi rõ họ tên)}
\end{tabular}
\end{flushright}

\end{document}
